% SPDX-License-Identifier: Apache-2.0
% covers: AxiPer
\datasheet{AxiPer}{The peripheral side of an AXI4 link: two address channels merged into one stream of requests.}

\dsfacts{\code{lib/parts/src/bus/axi.rs}}{\code{txhdl\_parts::bus::axi::AxiPer}}%
{\code{//docs:axi}, \code{//docs:vreteno}, \code{//docs:noc}}

\subsection*{Function}

\code{AxiPer} is the tracker at the peripheral end of an AXI4 link. It
merges \code{aw} and \code{ar} into one stream of \code{PerReq} on
\code{req}, each with a bit that says whether it is a read. It passes
the write beats from \code{w} to \code{wd}. It turns the client's
\code{Answer} on \code{ans} into \code{B} beats on \code{b}, and
passes the client's read beats from \code{rb} to \code{r}.

It allocates no identifiers. A response leaves with the identifier the
client gave it. Its ports are named types, \code{PerIn} and
\code{PerOut}, so a design may supply a unit of its own instead.
\code{//docs:axi} notes that a peripheral may also read the five AXI
channels directly and leave \code{AxiPer} out.

\subsection*{Parameters}

\begin{center}\footnotesize
\begin{tabular}{@{}lp{0.7\textwidth}@{}}
\toprule
Parameter & Meaning \\
\midrule
\code{A} & The address width, in bits. \\
\code{D} & The data width, in bits. \\
\code{S} & The write strobe width, one bit per byte lane, so \code{D / 8}. \\
\code{I} & The identifier width, in bits. \\
\bottomrule
\end{tabular}
\end{center}

There is no \code{NIDS}, since this end hands out no identifiers. The
tables below use \code{AxiPer<32, 32, 4, 2>}, the tracker in front of
Vreteno's data memory and timer.

\dstables{AxiPer}

\subsection*{Behaviour}

\begin{itemize}
\item At most one request leaves per cycle, and only when \code{req}
has room.
\item \code{rr} says which address channel goes first: zero for
\code{ar}. When both offer, the one that goes first is taken. When
only one offers, that one is taken.
\item Taking a read sets \code{rr} to one, and taking a write sets it
to zero. So two channels that both stream take turns, and neither
starves the other.
\item The request has \code{read} high when it came from \code{ar}.
Its \code{id}, \code{addr}, \code{len}, \code{size}, \code{burst},
\code{lock}, \code{cache}, \code{prot}, \code{qos} and \code{region}
are the address phase's.
\item A write beat passes from \code{w} to \code{wd} unchanged when
\code{wd} has room, independently of the requests.
\item An answer on \code{ans} becomes a \code{B} of the same \code{id}
and \code{resp} when \code{b} has room. A read beat passes from
\code{rb} to \code{r} unchanged when \code{r} has room.
\end{itemize}

\subsection*{Verification}

\begin{itemize}
\item Four tests in \code{bus::axi} run \code{AxiHost<16, 32, 4, 2, 4>}
and \code{AxiPer<16, 32, 4, 2>} with a client on each end:
\begin{itemize}
\item \code{a\_read\_reads\_what\_the\_write\_wrote};
\item \code{answers\_come\_back\_out\_of\_order};
\item \code{a\_transaction\_is\_answered\_by\_another\_process};
\item \code{every\_burst\_is\_answered\_and\_reads\_agree\_with\_writes}.
\end{itemize}
\item \code{bus::axi::fuzz} runs the two trackers with a relay that
stalls at random in every channel and a seeded client and peripheral
on the ends, and checks every burst answered once, reads against
writes, \code{last} on the final beat, and no deadlock, on 24 seeds;
\code{//lib/parts:axi\_fuzz\_test}, manual, runs more.
\item The tests of \code{bus::axi::sim} put the two trackers in front
of the simulated RAM:
\begin{itemize}
\item \code{a\_burst\_reads\_what\_a\_burst\_wrote};
\item \code{a\_strobe\_covers\_the\_lanes\_it\_says};
\item \code{a\_fixed\_burst\_stays\_on\_one\_word};
\item \code{an\_address\_past\_the\_end\_is\_answered};
\item \code{an\_image\_is\_loaded\_and\_read\_back}.
\end{itemize}
All of these run in \code{//lib/parts:parts\_test}.
\item \code{ex\_axi} has three bursts answered out of order. The build
simulates the netlist of \code{AxiPer<16, 32, 4, 2>} against that run
under nvc and Verilator:
\begin{itemize}
\item \code{//docs:axi\_sim\_axi\_per\_axi\_per\_tb\_test};
\item \code{//docs:axi\_vsim\_axi\_per\_test}.
\end{itemize}
\item The tests of the router, the Wishbone bridge and the network on
chip in \code{txhdl\_parts} run it in simulation. So do
\code{//ddr3:ddr3\_test}, \code{//gpu/razboj:razboj\_test},
\code{//soc:demo\_test} and \code{//cpu/vreteno:lockstep\_test}.
\end{itemize}

\subsection*{Used by}

The examples \code{ex\_axi}, \code{ex\_axi\_serve},
\code{ex\_axi\_router}, \code{ex\_axi\_reg} and \code{ex\_axi\_wb}.
Vreteno's run in \code{cpu/vreteno/src/run.rs} and
\code{cpu/vreteno/src/main.rs}, in front of the data memory and the
timer. The board's design, \code{Board}, as its fields \code{pdmem},
\code{ptimer} and \code{pddr3}. The four-corner system in
\code{soc/src/lib.rs}, and Razboj's simulation in
\code{gpu/razboj/src/sim.rs}.

\subsection*{Limits}

It does not count beats, look at \code{last}, or match write beats to
requests. A client does that: \code{Per::accept} gathers a write's
beats after its request. It does not check that an answer's identifier
belongs to an open request. It passes \code{lock}, \code{cache},
\code{qos} and \code{region} to the client and acts on none of them.
